% =====================================================================
%  chinese-review-booklet.tex  —  中文「期末复习手册 / 知识点速查」模板
%  逆向自 Router/毛概_期末知识点页码定位.pdf 的视觉系统：
%    青(#2F5C6E) + 玫红(#C0504D) 双色 · fancyhdr 页眉 · tcolorbox 圆角框
%    青底白字题号标签条 · navy 表头 + 斑马纹速查表
%  引擎：XeLaTeX / Tectonic（ctex + xeCJK，中文 fandol 字体自动下载）
%  编译：scripts/compile_tectonic.sh chinese-review-booklet.tex --preview
% =====================================================================
\documentclass[11pt]{ctexart}

% --- table 选项必须在 tcolorbox(其内部 \usepackage{xcolor}) 之前注入 ---
\PassOptionsToPackage{table,svgnames}{xcolor}
\usepackage[a4paper,margin=2.4cm,top=2.3cm,headsep=4mm]{geometry}
\usepackage{xcolor}
\usepackage[most]{tcolorbox}
\usepackage{fancyhdr}
\usepackage{enumitem}
\usepackage{array}
\usepackage{ragged2e}
\usepackage{pifont}   % \ding{72} 实心星，规避 Latin 字体缺 ★ 字形

% ===================== 配色（design tokens） =====================
\definecolor{rvteal}{HTML}{2F5C6E}   % 主色：表头 / 提示框 / 题号条
\definecolor{rvrose}{HTML}{C0504D}   % 强调：定位框 / 正文外框 / 项目符号
\definecolor{rvtealbg}{HTML}{EAF1F4} % 提示框浅底
\definecolor{rvrosebg}{HTML}{FBEEEE} % 定位框 / 斑马纹浅底
\definecolor{rvgray}{HTML}{6B6B6B}   % 副标题 / 页码灰
\definecolor{rvink}{HTML}{1A1A1A}

% ===================== 页眉（fancyhdr） =====================
\newcommand{\rvrunhead}{复习手册}   % 可用 \renewcommand 改运行页眉文字
\pagestyle{fancy}
\fancyhf{}
\renewcommand{\headrulewidth}{0pt}
\fancyhead[L]{\small\bfseries\color{rvteal}\rvrunhead}
\fancyhead[R]{\small\color{rvgray}第 \thepage\ 页}
\fancyhead[C]{\color{rvteal!45}\rule[-1mm]{0pt}{0pt}}
% 页眉下细线
\renewcommand{\headrule}{\vspace{1mm}{\color{rvteal!55}\rule{\headwidth}{0.4pt}}}

% ===================== tcolorbox 三种框 =====================
% 使用说明 / 提示（青）
\newtcolorbox{rvnote}{
  enhanced, breakable, colback=rvtealbg, colframe=rvteal,
  boxrule=0.7pt, arc=2.4mm, left=3.5mm, right=3.5mm, top=2.4mm, bottom=2.4mm,
  fontupper=\linespread{1.18}\selectfont}

% 教材定位 / 一行小标注（玫红，细）
\newtcolorbox{rvlocator}{
  enhanced, breakable, colback=rvrosebg, colframe=rvrose,
  boxrule=0.7pt, arc=2mm, left=3mm, right=3mm, top=1.5mm, bottom=1.5mm}

% 题目正文外框（玫红描边，白底）
\newtcolorbox{rvbody}{
  enhanced, breakable, colback=white, colframe=rvrose!85,
  boxrule=0.6pt, arc=2.4mm, left=3.6mm, right=3.6mm, top=2.6mm, bottom=2.2mm,
  fontupper=\linespread{1.18}\selectfont}

% ===================== 题号标签条（青底白字圆角） =====================
\newtcbox{\rvtab}{on line, enhanced, nobeforeafter,
  colback=rvteal, colframe=rvteal, boxrule=0pt, arc=1.6mm,
  left=2.6mm, right=2.6mm, top=1.1mm, bottom=1.1mm,
  fontupper=\bfseries\color{white}}
% \rvq{第 1 题}{标题}
\newcommand{\rvq}[2]{\medskip\par\noindent\rvtab{#1\hspace{1.1em}#2}\par\smallskip}

% ===================== 标题块 / 列表 / ★ =====================
\newcommand{\rvtitle}[2]{%
  \begin{center}
    {\LARGE\bfseries\color{rvink} #1\par}\vspace{2.2mm}
    {\large\color{rvgray} #2\par}
  \end{center}\vspace{1mm}}
\newcommand{\rvsectitle}[1]{\begin{center}\large\bfseries\color{rvteal} #1\end{center}\vspace{0.5mm}}
\newcommand{\impstar}{\textcolor{rvrose}{\ding{72}}}
\newcommand{\imp}{\textcolor{rvrose}{\ding{72}\,\bfseries 重要}}
\setlist[itemize]{leftmargin=1.7em, itemsep=2.5pt, topsep=2pt, parsep=0pt,
  label=\textcolor{rvrose}{\small\textbullet}}

% 表头单元格快捷
\newcommand{\rvth}[1]{\textcolor{white}{\bfseries #1}}

\begin{document}
\renewcommand{\rvrunhead}{示例复习手册 · 知识点速查}

% ---------- 标题 ----------
\rvtitle{XXX 课程·期末核心复习手册}{（202X 版）· 知识点速查与要点提纲}

% ---------- 使用说明（青框） ----------
\begin{rvnote}
\textbf{使用说明：} 本资料把每个复习知识点当作一道题，给出\textbf{定位}与\textbf{可背诵要点}。
标注 \imp{} 为划定的重点题。要点依教材原文与原结构整理。圆角框、双色版式逆向自「毛概」样式。
\end{rvnote}

\vspace{2mm}
% ---------- 速查表（navy 表头 + 斑马纹） ----------
\rvsectitle{知识点速查表}
{\renewcommand{\arraystretch}{1.35}%
\rowcolors{2}{rvrosebg}{white}%
\begin{center}
\begin{tabular}{>{\centering\arraybackslash}m{2.6em} >{\RaggedRight}m{7.2cm}
                >{\centering\arraybackslash}m{2.4cm} >{\centering\arraybackslash}m{2.4cm}}
\rowcolor{rvteal}\rvth{题号} & \rvth{知识点} & \rvth{印刷页码} & \rvth{速查页} \\
1 & 第一个核心知识点的内涵       & 4--6   & 16--18 \\
2 & 第二个知识点及其相互关系     & 10--12 & 22--24 \\
3 & 第三个知识点（含子结构）\imp & 26--32 & 38--44 \\
4 & 第四个知识点的历史地位       & 32--37 & 44--49 \\
\end{tabular}
\end{center}}

% ---------- 题目（青标签条 + 玫红定位 + 玫红正文框） ----------
\rvq{第 1 题}{第一个核心知识点的内涵}
\begin{rvlocator}\textbf{教材定位：}印刷页 4--6（速查第 16--18 页）\end{rvlocator}
\begin{rvbody}
\textbf{书中标题：}一、第一个核心知识点的内涵\quad\textbf{复习要点：}
\begin{itemize}
  \item \textbf{提出必要性：}写明为什么需要这一知识点——一句话点出动因与背景。
  \item \textbf{总定义：}给出可逐字背诵的标准定义，保留教材用语。
  \item \textbf{内涵分层：}第一层（…）、第二层（…）、第三层（…），逐条展开。
  \item \textbf{把握要点 \imp：}考试高频的辨析点，正反对照写清。
\end{itemize}
\end{rvbody}

\rvq{第 2 题}{第二个知识点及其相互关系}
\begin{rvlocator}\textbf{教材定位：}印刷页 10--12（速查第 22--24 页）\end{rvlocator}
\begin{rvbody}
\textbf{书中标题：}二、第二个知识点及其相互关系\quad\textbf{复习要点：}
\begin{itemize}
  \item \textbf{构成：}列出该知识点的组成部分与边界。
  \item \textbf{相互关系：}与第 1 题的联系——一脉相承 / 与时俱进，分两面写。
  \item \textbf{共同地位：}收束一句总评。
\end{itemize}
\end{rvbody}

\end{document}
